\section{Layered scoring and scientific completion}
\label{sec:si-scoring}

\subsection{Composite and reporting layers}
For a formally organized run, the reported composite combines graph structure $G$, physical self-consistency $C$, and accuracy $A$:
\begin{equation}
S=0.1G+0.3C+0.6A.
\label{eq:si-composite}
\end{equation}
The three components are retained separately. Execution status, original graph delivery, completion of requested steps, graded physical evidence, and complete scientific acceptance are also reported separately. A completed process or a partial score does not substitute for the evidence required by a particular scientific claim.

\subsection{Graph structure and chemical correspondence}
The graph component measures coverage of the reference mechanism by the delivered mechanism. Transition-state nodes are folded into their elementary-step edges for graph comparison. Species are matched chemically and reaction steps are compared through their endpoints and atom-mapped bond changes. RDKit supplies chemical representations and NetworkX supplies graph representation and matching.\citep{ref13,networkx2008} Missing reference elements, incorrect chemical identities, or mismatched accessibility declarations contribute edit costs. Additional delivered nodes or edges are not penalized merely for being additional in this component.

The tier-specific graph comparison ignores stereochemistry for Easy and Normal tasks and retains it for Hard tasks. Bound native evidence determines the evidence association of a delivered species. The graph score is not multiplied by the binary self-consistency pass flag: graph coverage and physical self-consistency remain distinct components.

\subsection{Graded physical self-consistency}
The self-consistency assessment operates on delivered species and steps, independently of numerical agreement with Gold. A transition-state step receives graded credit for the evidence stages in Table~\ref{tab:si-self-consistency}. A complete verified chain receives 1.0. When necessary, evaluation searches the same run's native artifacts for the corresponding chain; recovery of a file association retains the same physical acceptance criteria.

\begin{table}[htbp]
\centering
\caption{Graded self-consistency contribution for an accessible transition-state step. The endpoint contribution is 0.25 for one verified side and 0.50 for both sides.}
\label{tab:si-self-consistency}
\small
\begin{tabular}{p{0.65\textwidth}r}
\toprule
Evidence stage & Contribution \\
\midrule
Verified first-order saddle: exactly one imaginary frequency & 0.35 \\
IRC starting geometry consistent with the selected saddle & 0.15 \\
Native outer-frame--optimization--frequency lineage to both declared endpoint minima & 0.50 \\
\bottomrule
\end{tabular}
\end{table}

For graphs containing reaction edges, $C$ is the mean of the delivered edge scores. An edge without a transition-state node requires both endpoint species to be covered. For a species-only graph, $C$ is the mean species-coverage score. Explicit inaccessible targets are exempt from requirements for physically nonexistent stationary points or paths and receive the corresponding self-consistency credit.

The complete self-consistency pass is a separate all-required-elements condition. It checks the required species and steps, including minimum evidence and the full transition-state/IRC chain. Consequently, a full graded component and a binary scientific-acceptance record are retained as separate outputs. Counts of $C=1$ in the main text and supplementary score tables refer to full component scores; they are not renamed as counts of fully accepted complete mechanisms. Fractional scores identify the attained stages of physical support.

\subsection{Electronic barriers and numerical accuracy}
For a numerical reference step, the declared barrier is checked against an energy difference derived from its native calculation evidence. An omitted edge barrier can be obtained from the delivered transition-state and starting-species energies when those quantities and their native records support the conversion. The energy difference is converted from hartree to kcal~mol$^{-1}$ using the common conversion factor 627.509474.

For a supported numerical value $\widehat{\Delta E}_e^{\ddagger}$ and reference $\Delta E_{e,\mathrm{Gold}}^{\ddagger}$, the underlying Gaussian score is
\begin{equation}
q_e=\exp\left[-\frac{\left(\widehat{\Delta E}_e^{\ddagger}-
\Delta E_{e,\mathrm{Gold}}^{\ddagger}\right)^2}{2\sigma_e^2}\right].
\label{eq:si-gaussian}
\end{equation}
The released task rubrics use base widths of 3.0~kcal~mol$^{-1}$ for Easy and Impossible tasks and 2.0~kcal~mol$^{-1}$ for Normal and Hard tasks. For a reference barrier at or below 10.0~kcal~mol$^{-1}$, the reference-comparison width is the larger of the base width and 5.0~kcal~mol$^{-1}$. This small-barrier width applies to comparison with Gold; the declaration-to-log consistency check retains the task's base width.

All reference edges remain in the accuracy denominator, with their recorded weights. A numerical reference edge lacking a matching delivered step, a declared barrier, or a log-derived value receives zero for that edge. Other evidenced steps retain their contributions. Numerical deviations on supported reference edges are aggregated with a maximum penalty fraction of 0.3; missing-step/evidence penalties and incorrect negative conclusions are not protected by this numerical-deviation cap. More explicitly, for total reference-edge weight $W$, supported numerically graded edges $D$, and missing or incorrectly delivered targets $U$, the aggregation takes the form
\begin{equation}
A=1-\min\!\left(\frac{\sum_{e\in D}w_e(1-q_e)}{W},\,0.3\right)
-\frac{\sum_{e\in U}w_e}{W}.
\label{eq:si-accuracy-aggregation}
\end{equation}
Correct negative targets and accessible edges intentionally lacking a numerical reference contribute their prescribed full-credit terms through $W$. The latter are not interpreted as numerical barrier predictions.

The evaluator considers the legal transition-state/reference-energy combinations present in the delivered evidence. When those records contain a combination agreeing with Gold but the reported value uses a different supported conformer or energy-reference definition, the documented answer-definition treatment assigns an edge score of 0.9 in place of $q_e$. A declared barrier that is unsupported by any legal evidence combination, with no reference-consistent alternative in the delivered records, triggers the integrity short circuit: the composite retains only $0.1G$, and the self-consistency and accuracy contributions are zeroed. An omitted step is handled by its missing contribution rather than by this inconsistent-declaration rule.

\subsection{Qualitative selectivity and property targets}
HA1, HA6, HB1, and HF1 use a qualitative favored-channel target in the accuracy slot. The delivered mechanism must identify one favored edge corresponding to the reference channel or its chemically matched product. Product identity supplies the chemical correspondence when the delivered channel label differs from the internal reference label. The numerical channel spread is not scored as a quantitative selectivity prediction in these tasks. HL1 instead retains its independently referenced channel barriers and does not use this favored-channel substitution.

When a task specifies species properties, its accuracy assessment compares the declared properties with the corresponding reference targets. Numerical property widths are 0.02~\AA\ for bond lengths, 50~cm$^{-1}$ for frequency quantities, and 2.0~kcal~mol$^{-1}$ for energy differences; the specific keyed property determines the width. Specified state-ordering text is compared after case and whitespace normalization. Multiple specified properties contribute their mean score.

\subsection{Impossible tasks and final scientific claims}
Impossible-task assessment is based on the final delivered target. Correctly identifying that a requested target is inaccessible or outside the allowed method's capability earns the corresponding negative-target credit. An exploratory calculation followed by a correct negative conclusion is not penalized for having been attempted. An incorrect positive accessibility claim, a purportedly trustworthy barrier/property value for the negative target, or an absent required negative conclusion loses the corresponding credit. Negative endpoint species are treated consistently with their associated negative step. A species-only negative target is assessed without introducing a process edge.

The result must be interpreted at the identity of the requested target. A valid result for another electronic state, species, or pathway does not automatically answer the stated question. The per-task tables and case analyses retain this target association alongside the numerical score.

\subsection{Original delivery, organized evaluation, and score reproduction}
Original delivery completeness records whether the run generated its mechanism graph during execution. Organized evaluation records the graph obtained by the formal offline procedure in Section~\ref{sec:si-protocol}. The publication component table preserves each run's original or organized score record and field provenance; existing released scores are reproduced rather than replaced by scores from a later software environment.

The retained reproduction procedure joins the released run table to the corresponding recorded components and reconstructs the aggregate tables and figure data. Its checks cover 200 unique task/configuration records, 600 nonmissing component values, and consistency of the weighted components with the 200 released totals. The maximum absolute difference from the published four-decimal composites is $4.9711\times10^{-5}$. The scoring formulas above explain the layered evaluation; regenerating publication tables requires the retained component records, while independently reassessing native scientific evidence additionally uses the associated run artifacts and evaluator revision.
